\RequirePackage{fix-cm}
\documentclass[9pt,letterpaper]{extarticle}
\usepackage[margin=0.38in,top=0.35in,bottom=0.53in,footskip=18pt]{geometry}
\usepackage{lmodern}\usepackage{fontspec,amsmath,amssymb,array,tabularx,fancyhdr,xcolor,hyperref,enumitem}
\setmainfont{Latin Modern Sans}\setsansfont{Latin Modern Sans}
\definecolor{Accent}{HTML}{145E68}\definecolor{Muted}{HTML}{4D555B}
\hypersetup{hidelinks,pdftitle={Water Quality B 2027 — Exam reference},pdfauthor={Greg Hogan and the soso contributors}}
\input{@@NOTICE@@}
\pagestyle{fancy}\fancyhf{}\renewcommand{\headrulewidth}{0pt}
\fancyfoot[L]{\sffamily\fontsize{6.2}{7}\selectfont\sosocredit}
\renewcommand{\footrulewidth}{0.3pt}
\fancyfoot[R]{\sffamily\fontsize{8}{9}\selectfont\thepage\,/\,2}
\setlength{\parindent}{0pt}\setlength{\parskip}{2pt}\setlength{\emergencystretch}{12pt}
\setlength{\tabcolsep}{2pt}\renewcommand{\arraystretch}{1.1}
\newcommand{\head}[2]{{\sffamily\bfseries\fontsize{13}{15}\selectfont WATER QUALITY B}\hfill{\sffamily\fontsize{8.5}{10}\selectfont 2027\quad|\quad EXAM REFERENCE}\par\vspace{2pt}\hrule height0.6pt\vspace{4pt}{\sffamily\bfseries\fontsize{10.3}{12}\selectfont #1}\par{\sffamily\fontsize{7.4}{9}\selectfont #2}\par\vspace{3pt}}
\newcommand{\sect}[1]{\par\vspace{5pt}{\sffamily\bfseries\color{Accent}\fontsize{9.2}{11}\selectfont #1}\par\vspace{2pt}{\color{Accent}\hrule height0.35pt}\vspace{3pt}}
\newcommand{\term}[2]{\par\vspace{0.7pt}\textbf{#1}\enspace #2\par}
\newenvironment{cols}{\begin{minipage}[t]{0.484\textwidth}\raggedright}{\end{minipage}}
\newcommand{\nextcol}{\end{minipage}\hfill\begin{minipage}[t]{0.484\textwidth}\raggedright}
\begin{document}\fontsize{9}{10.4}\selectfont\setlength{\abovedisplayskip}{3pt}\setlength{\belowdisplayskip}{3pt}\setlength{\abovedisplayshortskip}{2pt}\setlength{\belowdisplayshortskip}{2pt}
\head{1 · Identify organisms; connect them to the ecosystem}{Print both pages on opposite sides of ONE US Letter sheet. No affixed labels. Common names are required for identification.}
\begin{cols}
\sect{Reef identification: shape before color}
Use body plan, fins, appendages, surface texture and habitat together. Color varies with age, sex, light and species. Size thresholds require scale; do not infer length from a photograph alone.
\term{Banded coral shrimp}{Red/white bands, long white antennae and prominent front claws; cleaner and scavenger. Eggs carried by female; planktonic larvae.}
\term{Butterflyfish}{Deep, thin body, small projecting mouth; often eye bar or false eyespot. Diet varies: coral polyps, small invertebrates or plankton. Coral-feeders may track coral condition.}
\term{Crown-of-thorns starfish}{Many arms with dense long spines; coral predator. Outbreaks can remove living coral tissue; do not confuse with an ordinary five-armed star.}
\term{Fleshy algae}{Soft fronds, blades or turf; red, green or brown groups. Producers; excess cover may reflect nutrient enrichment, lost grazers or newly available substrate, not one unique cause.}
\term{Grouper $>30$ cm}{Robust fish, large mouth, strong tail; predator. Larger individuals may be reduced by fishing. Includes a wide range of patterns.}
\term{Hard/stony corals}{Colonies or solitary polyps with rigid calcium-carbonate skeletons. Many reef builders host photosynthetic symbionts and also capture food.}
\term{Spiny lobsters}{Long stiff antennae, spiny body, \emph{no huge paired front claws}; Palinuridae, not clawed lobsters. Often nocturnal foragers; long planktonic larval stage.}
\term{Long-spined black urchins}{Round test with very long thin spines; algae grazers. Grazing can open space for coral recruitment. Distinguish from stout-spined pencil urchins.}
\term{Moray eel}{Elongate fish, continuous dorsal fin and no pectoral fins; crevice predator. Mouth opening pumps water over gills and is not by itself aggression.}
\term{Parrotfish $>20$ cm}{Fused beak-like teeth; many scrape algae and material from reef surfaces. Grazing and bioerosion; colors often change through life. Corrected group: Scarinae.}
\term{Pencil urchins}{Thick, blunt, pencil-like spines; \emph{Eucidaris}. Omnivorous grazing differs from the fine-spined black urchins.}
\term{Sea cucumbers}{Soft elongated echinoderms with feeding tentacles; many process sediment or suspension-feed. Nutrient recycling; larvae develop into bottom-dwelling adults.}
\term{Snappers}{Fish with relatively large mouth and canine teeth; predators of fish/invertebrates. Use body/fin shape, not red color alone.}
\term{Demosponges}{Porous bodies with water outlets, no coral polyps; filter feeders. Water moves through canals; shapes include tubes, barrels and encrusting forms.}
\term{Sweetlips}{Fleshy lips; many juveniles have striking spots and different adult patterns. Often nocturnal benthic-invertebrate feeders; daytime sheltering.}
\term{Triton}{Large spiral snail with elongated shell opening; predator of echinoderms, including crown-of-thorns. A shell alone may obscure the living-animal features.}
\sect{Treatment and contaminants}
\term{Drinking water}{Coagulation/flocculation, settling and filtration remove particles; disinfection targets microbes. Dissolved contaminants may need additional treatment.}
\term{Wastewater}{Primary settling removes solids; secondary biological treatment reduces organic load. Nutrient removal/disinfection are additional steps.}
\term{HABs / emerging pollutants}{Not all blooms are toxic; harms include shading and oxygen loss. Pharmaceuticals/PFAS can escape treatment. Detection alone does not quantify risk.}
\nextcol
\sect{Regional targets in the event list}
\textbf{Indo-Pacific group}\par
\term{Humpback grouper / barramundi cod}{High arched back, small head, dark spots on pale body; \emph{Cromileptes altivelis}. Predator; juvenile appearance differs in proportions.}
\term{Bumphead parrotfish}{Large blunt forehead and exposed beak; \emph{Bolbometopon muricatum}. Often schools; grazes/erodes reef material.}
\term{Giant clams}{Heavy, fluted paired shells and colorful exposed mantle; \emph{Tridacna}. Filter feeding plus photosynthetic symbionts; do not identify from mantle color alone.}
\term{Humphead wrasse}{Large adult forehead hump, thick lips; \emph{Cheilinus undulatus}. Juveniles lack the pronounced hump. Predator of invertebrates and fish.}
\textbf{Atlantic group}\par
\term{Flamingo tongue snail}{Patterned living mantle covers the shell; withdrawn mantle reveals a largely unpatterned shell. \emph{Cyphoma gibbosum} grazes soft corals.}
\term{Sea fans / gorgonians}{Branching flexible colonies, often flat mesh-like fans; octocorals have eightfold polyp symmetry. Suspension feeding; many host symbionts.}
\term{Nassau grouper}{Five dark body bars, dark saddle near tail base; \emph{Epinephelus striatus}. Predator; spawning aggregations are vulnerable to fishing.}
\sect{Reefs: evidence, not a health label}
\term{Bleached vs recently killed}{Bleaching: living coral loses symbionts/pigments; white skeleton shows through tissue. Death: tissue is absent; algae may colonize skeleton. Bleaching can recover or lead to death.}
\term{Cause chains}{Heat stress $\to$ bleaching; sediment $\to$ shading/smothering; lost grazers + nutrients $\to$ more algae; $\mathrm{CO_2}$ uptake $\to$ lower pH/carbonate availability. Distinguish mechanism from correlation.}
\term{Life cycles}{Many reef animals release gametes or brood eggs, then have dispersing larvae and a settled juvenile stage. Recruitment depends on suitable substrate, food and survival; adult abundance is not just reproduction rate.}
\sect{Marine and estuary ecology}
\term{Estuary}{Freshwater and seawater mix; tides, river flow and density create variable gradients. Nursery habitat, nutrient processing and shoreline protection; salinity stress selects tolerant organisms.}
\term{Watershed}{Land draining to a shared outlet. Upstream fertilizer, erosion, sewage and altered flow affect downstream estuaries even without a nearby discharge.}
\term{Food webs}{Producers fix energy; consumers transfer it; decomposers recycle nutrients. Energy dissipates as heat, while matter cycles. A trophic-transfer percentage is not universal.}
\term{Interactions}{Predation, competition, mutualism, parasitism and facilitation alter populations. Limiting resources and carrying capacity vary; removing a predator can cause a trophic cascade.}
\term{Carbon/nitrogen}{Photosynthesis consumes $\mathrm{CO_2}$; respiration releases it. Nitrogen fixation supplies usable N; nitrification consumes oxygen; denitrification can remove nitrate as gas in low-oxygen conditions.}
\term{Garbage patches}{North/South Pacific, North/South Atlantic and Indian Ocean gyres concentrate dispersed debris, not solid islands. Entanglement, ingestion and microplastic exposure differ from chemical toxicity.}
\end{cols}
\newpage
\head{2 · Interpret measurements; calibrate the device}{Ecology 30\% · organism identification 30\% · monitoring/data 30\% · salinometer 10\%. Check current rules and clarifications.}
\begin{cols}
\sect{Read the measurement in context}
\term{Salinity}{Salts affect density, osmosis and species distribution. Evaporation concentrates; freshwater dilutes. Conductivity responds to ions \emph{and temperature}.}
\term{Temperature}{Warmer water generally holds less oxygen at saturation and changes metabolic rates. Compare like depths/times; stratification can isolate bottom water.}
\term{Dissolved oxygen (DO)}{mg/L or \% saturation. Photosynthesis/aeration add oxygen; respiration/decay consume it. Daytime surface DO can mask low pre-dawn/bottom DO.}
\term{Biochemical oxygen demand (BOD)}{Oxygen consumed by microbes during a specified incubation; high BOD indicates oxygen-demanding material. High BOD is not high available DO. BOD$_5$ uses a five-day protocol.}
\term{pH}{ $\mathrm{pH}=-\log_{10}a_{\mathrm{H^+}}$. A one-unit fall means tenfold greater hydrogen-ion activity. Ocean acidification means falling pH, not necessarily pH below 7.}
\term{Aragonite saturation $\Omega$}{Ion activity product divided by solubility product. $\Omega<1$: dissolution favored; $\Omega>1$: supersaturated, but does not guarantee biological calcification or reef health.}
\term{Nitrate / phosphate}{Nutrients can limit growth; enrichment can fuel blooms. A low value during a bloom may reflect uptake, not absence of enrichment. Check units and whether concentration is reported ``as N/P'' or as the whole ion.}
\term{Turbidity}{Light scattering, often NTU; caused by suspended particles/plankton. Reduces light; deposited sediment can smother habitat. It is not the same as dissolved solids or an exact sediment concentration.}
\term{Total solids}{Dissolved + suspended material under a specified test method. Conductivity mainly tracks ions; turbidity tracks scattering material.}
\term{Fecal coliform}{Indicator bacteria suggesting fecal contamination; not a count of all pathogens or proof of one specific source. Units and sampling method matter.}
\sect{How measurements are made}
Calibrate pH probes with buffers; DO probes need calibration and temperature/salinity correction. BOD compares initial/final DO after incubation with dilution/blank corrections. Turbidity uses a calibrated light-scattering meter; nutrients use standards and colorimetry. Total solids: dry a known volume and weigh residue. Fecal indicators: culture/count under a specified protocol, not visual inspection.
\sect{Data interpretation checklist}
Check axes, units, uncertainty, time/depth, tide and rainfall. Describe patterns before mechanisms: correlation is not causation. Test alternative explanations with extra measurements; replicate and use blanks/standards.
\sect{Common coupled patterns}
Nutrients $\uparrow$ $\to$ growth $\uparrow$ $\to$ decay/BOD $\uparrow$ $\to$ DO $\downarrow$; daytime photosynthesis may temporarily reverse DO/pH patterns.\par
Rain can reduce estuary salinity while increasing turbidity, nutrients and fecal indicators. Density gradients may reduce vertical oxygen renewal.\par
A low-DO measurement alone cannot distinguish warming, poor mixing, high respiration or contaminated runoff.
\nextcol
\sect{Hydrometer: what the float measures}
A floating object balances weight and buoyancy:
\[m g=\rho_{\rm solution}gV_{\rm displaced}.\]
Denser solution requires \emph{less displaced volume}; the hydrometer floats higher. Equilibrium buoyancy still equals its weight.
\term{Stable design}{A sealed narrow stem gives a readable scale; ballast low in the body gives stability. Nothing should leak, dissolve, touch the beaker wall or rest on its bottom. Remove trapped bubbles.}
\term{Calibration}{Use several known concentrations across the range; repeat readings at a controlled temperature. Plot concentration against stem reading; do not assume a globally linear scale. Check with a blind standard.}
\term{Read consistently}{Let motion settle. Use eye level and the same meniscus convention used in calibration. Record units, temperature, repeatability and reading direction; do not extrapolate beyond the calibrated range.}
\sect{Concentration: name the denominator}
\[\%\,(m/v)=\frac{m_{\rm solute}\ (\mathrm g)}{V_{\rm final\ solution}\ (\mathrm{mL})}\times100.\]
For 250 mL of 4.0\% $m/v$: weigh 10.0 g salt, dissolve in less than 250 mL water, then bring the \emph{final solution volume} to 250 mL. Adding 10 g to 250 mL of water is not the same procedure.\par
Mass percent uses mass of \emph{solution}; practical salinity is conductivity-based. Do not freely interchange $\%m/v$, $\%m/m$, g/kg, ppt and practical salinity. $1\%m/v=10$ g/L by definition.
\sect{Interpolation and uncertainty}
Between nearby calibration points $(x_1,C_1)$, $(x_2,C_2)$:
\[C\approx C_1+(C_2-C_1)\frac{x-x_1}{x_2-x_1}.\]
Let $x$ be exposed-stem length above the meniscus: greater $x$ means higher concentration. If 3.0\% and 5.0\% standards read 24 and 36 mm, a 30 mm reading estimates 4.0\%. Validate locally; curvature, temperature, salt composition and drift can bias it.\par
\textbf{Precision} is repeatability; \textbf{accuracy} is closeness to a reliable target. Repeated results can share a calibration bias. A tenth-percent display does not establish tenth-percent accuracy.
\sect{Calculate and compare evidence}
\term{Oxygen saturation}{\(\%\,\text{saturation}=100\times\mathrm{DO}_{\rm measured}/\mathrm{DO}_{\rm saturation}\), with both concentrations in the same units. Use the saturation value for the sample's temperature, salinity and pressure. A value above 100\% can follow strong photosynthesis; it does not prove lasting good conditions.}
\term{pH change}{A drop from pH 8.2 to 7.9 multiplies hydrogen-ion activity by \(10^{8.2-7.9}\approx2\). This is not a 0.3\% change; pH is logarithmic.}
\term{Abundance and effort}{Density = count / sampled area; relative abundance = count of a group / total count. Compare equal sampling effort, habitat and detectability. Species richness counts species, not individuals; evenness describes how equally individuals are distributed among species.}
\term{Change versus percentage points}{Percent change = \(100(\text{new}-\text{old})/\text{old}\), for a nonzero baseline. Coral cover falling from 40\% to 30\% is a 10-percentage-point fall, but a 25\% relative decline. State which measure you used.}
\vspace{3pt}{\sffamily\fontsize{6.8}{8}\selectfont Basis: 2027 B rules pp. B59--60; national correction 6 Nov 2026; original synthesis of NOAA coral/estuary and USGS water-science resources. Rules checked 17 Nov 2026. Full sources and explanations are in the study/build guides. \url{soinc.org/events/rules-corrections}}
\end{cols}
\end{document}
